\chapter{Literature Review}
\setcounter{equation}{0}
\sloppy

\section{XGBoost-Based URL Phishing Detection Method With Cross-Dataset Validation, 2026}
The paper by Milosz Misiek and Tomasz Hyla presents a machine-learning approach
for detecting phishing URLs using XGBoost. The study addresses the difficulty of
identifying malicious URLs that resemble legitimate web addresses and the
challenge of maintaining detection accuracy when the model is evaluated on
datasets different from those used during training. The authors emphasize that
URL characteristics vary across datasets, which can affect the generalization
ability of phishing-detection models.

The proposed method uses Python-based feature extraction to derive URL
characteristics, including lexical features, DNS and network information,
security-related attributes, domain reputation, and redirect behavior. The study
extracts more than 90 candidate features and applies feature selection to identify
characteristics that remain useful across datasets. An XGBoost classifier is then
evaluated using three datasets and compared with neural-network approaches. The
paper reports 90.7\% accuracy on its custom dataset and 77.8\% average accuracy
across three independent test datasets. These results highlight the importance of
feature selection and cross-dataset evaluation.

BrowseBlock adopts XGBoost for URL phishing classification and feature-based URL
analysis. The extracted URL characteristics are supplied to the trained classifier
to determine whether a destination appears legitimate or suspicious. The study's
emphasis on cross-dataset validation is relevant when evaluating BrowseBlock
against previously unseen URLs. For QR-code phishing, BrowseBlock adds a decoding
stage that extracts the embedded URL and sends it through the same URL-analysis
pipeline; QR decoding itself is not the principal method proposed in this paper.

\clearpage
\section{Classification of Phishing Email Using Random Forest Machine Learning Technique, 2014}
The paper by Andronicus A. Akinyelu and Aderemi O. Adewumi investigates the use
of machine learning to classify phishing emails. It addresses limitations of
traditional email-filtering methods, which may fail to recognize newly developed
phishing patterns because attackers frequently change the wording and structure of
their messages. The research aims to identify useful email characteristics that
distinguish phishing messages from legitimate messages while reducing the number
of features required for classification.

The proposed approach extracts phishing-related features from a dataset of 2,000
emails and converts them into feature vectors for machine-learning classification.
The authors use Random Forest and evaluate different feature combinations,
including a reduced set of eight prominent features. The paper reports 99.7\%
classification accuracy for its evaluated model. This result reflects the dataset
and evaluation procedure used in the study and should not be assumed to represent
the accuracy of every phishing-detection system.

BrowseBlock adopts Random Forest as the initial email-phishing screening technique.
The classifier evaluates available email features and identifies messages that
require further analysis. BrowseBlock extends this approach by examining URLs
extracted from email content, incorporating SPF, DKIM, and DMARC results when
available, and using RAG with an LLM for deeper contextual analysis. These
additional stages are part of BrowseBlock's proposed hybrid pipeline rather than
technologies introduced by the Random Forest paper.

\clearpage
\section{Detecting DOM-Sourced Cross-Site Scripting in Browser Extensions, 2017}
The paper by Jinkun Pan and Xiaoguang Mao investigates DOM-sourced Cross-Site
Scripting vulnerabilities in browser extensions. Extension scripts can interact
with webpage content and may possess privileges beyond those available to ordinary
webpage scripts. When attacker-controlled webpage data is processed unsafely by
an extension, it may cause malicious script execution or unauthorized access to
protected resources. The research focuses on identifying these vulnerabilities,
particularly in user scripts associated with the Greasemonkey browser extension.

The authors propose a hybrid analysis framework with two phases. The first phase
performs lightweight static analysis using text filtering and an Abstract Syntax
Tree parser to identify potentially vulnerable code. The second phase applies
dynamic symbolic execution together with a shadow-DOM component and
constraint-solving techniques to investigate candidate vulnerabilities. The paper
reports 58 previously unknown DOM-sourced XSS vulnerabilities in Greasemonkey
user scripts. The results demonstrate the usefulness of combining lightweight
source analysis with deeper investigation of potentially unsafe execution paths.

BrowseBlock adopts the concept of examining suspicious data flows between
untrusted webpage inputs and potentially dangerous JavaScript operations. Its XSS
module can inspect accessible scripts and identify indicators involving sources
such as URL parameters and sinks such as unsafe DOM-manipulation functions. The
paper provides a reference for more advanced static and dynamic analysis
techniques; BrowseBlock should claim AST analysis, symbolic execution, or
shadow-DOM analysis only if those techniques are actually implemented.

\clearpage
\section{Hidden Path: Understanding the Intermediary in Malicious Redirections, 2022}
The paper by Yuwei Zeng and co-authors examines how attackers use intermediary
websites and redirect chains to conceal malicious destinations. Malicious redirects
can make a user's navigation path difficult to understand by routing the browser
through several domains before reaching the final page. Attackers may use delays,
fabricated user interactions, or visitor characteristics to avoid detection by
conventional security checks. The study investigates mechanisms used in these
campaigns and the role of intermediate domains.

The researchers develop an iterative framework that uses a bipartite graph-based
method to discover domains potentially associated with malicious redirections from
real-world DNS traffic. They then use a Chrome-based crawler to visit suspicious
domains and reconstruct redirect chains from browser performance logs. The analysis
examines HTTP-, HTML-, and JavaScript-based navigation and identifies redirection
fluxing, in which randomness is introduced into the redirect process to make
malicious destinations harder to track.

BrowseBlock adopts the concepts of monitoring navigation changes, examining
suspicious redirect chains, and assessing the final destination using URL analysis.
The redirect module can evaluate the requested URL and subsequent destinations
where the required navigation information is available, and submit suspicious
destinations to the XGBoost URL detector. The paper's graph-based DNS mining and
large-scale crawling framework are reference techniques; they must not be
described as implemented BrowseBlock components unless separately developed.

\clearpage
\section{Implementation of Extension for Browser to Detect Vulnerable Elements on Web Pages and Avoid Clickjacking, 2016}
The paper by Dipti Pawade, Abhilasha Lahigude, Divya Reja, and Era Johri presents
ClickProtect, a browser extension designed to identify webpage elements vulnerable
to clickjacking. Clickjacking tricks a user into clicking an element different from
the one intended. The paper examines invisible iframe overlays and fake cursors
that conceal the actual target of a user interaction, with the objective of
detecting suspicious elements and warning the user before an unsafe interaction.

The proposed system consists of a revealing unit and an extenuation unit. The
revealing unit inspects webpage source code and presentation characteristics to
identify potentially hidden or misleading elements. The extenuation unit makes
suspicious elements visible or otherwise helps prevent the user from interacting
with a concealed target. This separation is relevant to BrowseBlock because
detection and user protection are separate responsibilities even when they are
presented in one browser workflow.

BrowseBlock adopts the concepts of inspecting accessible webpage elements for
suspicious visibility, overlay, and interaction indicators and warning the user
when findings are concerning. These checks contribute to its clickjacking
detection module, which can present supporting evidence through the extension
popup and dashboard. Since legitimate websites also use iframes and overlays,
the indicators should be combined with contextual checks rather than treating
every iframe or invisible element as proof of an attack.

\clearpage
\section{Analyzing In-Browser Cryptojacking, 2024}
The paper by Muhammad Saad and David Mohaisen investigates in-browser
cryptojacking, in which a webpage executes JavaScript that uses a visitor's
computing resources to mine cryptocurrency without appropriate consent. Unlike
conventional malware requiring a separate executable, in-browser cryptojacking
operates through scripts loaded by a webpage. It can consume CPU resources, drain
batteries, and affect normal system performance.

The researchers combine static and dynamic analysis. Static analysis extracts
code-related features, while dynamic analysis investigates mining effects on CPU
usage and battery consumption. The paper evaluates Logistic Regression, Linear
Discriminant Analysis, k-Nearest Neighbours, Support Vector Machines, and Random
Forest to distinguish cryptojacking scripts from benign and other malicious
JavaScript samples. It also examines browser communication and the economic
feasibility of cryptojacking as an alternative to online advertising.

BrowseBlock adopts static JavaScript inspection, suspicious-script feature
extraction, and machine-learning-based classification for cryptojacking detection.
Resource-usage indicators may also help when the required information is available,
but CPU consumption alone cannot prove cryptojacking. BrowseBlock's claims must
therefore reflect the script indicators and runtime information actually accessible
to its browser extension.

\clearpage
\section{Hidost: A Static Machine-Learning-Based Detector of Malicious Files, 2016}
The paper by Nedim Šrndic and Pavel Laskov introduces Hidost, a static
machine-learning system for detecting malicious files. It addresses malware hidden
inside document formats such as PDF and SWF. These files are commonly legitimate,
but their complex internal structures can be exploited through obfuscation,
encoding, or manipulation of document elements. The research aims to improve
detection without execution-based analysis.

Hidost combines information about a file's logical structure with its content to
generate features for machine-learning classification. It performs static analysis
without executing the file and was evaluated on datasets containing approximately
440,000 PDF files and 40,000 SWF files. Structural information characterizes files
and improves detection of malicious documents. The reported experiments showed
strong detection performance, particularly for malicious PDF files, compared with
the evaluated antivirus engines.

BrowseBlock adopts static file analysis, file-structure and content-feature
extraction, and machine-learning-based malicious-file classification. These
methods can assess supported downloads without executing them. The proposed
BrowseBlock file pipeline can combine static features with SHA-256 reputation
checks and a trained classifier, if those components are implemented. SHA-256
matching is an additional reputation mechanism, not the central technique
introduced by Hidost.

\clearpage
\section{When Push Comes to Ads: Measuring the Rise of (Malicious) Push Advertising, 2020}
The paper by Karthika Subramani and co-authors examines web push notifications
used to deliver advertisements, including malicious campaigns. Web push
notifications can reach users even when the original webpage is no longer in the
foreground. Although legitimate applications use notifications, attackers can
distribute deceptive advertisements and direct users to scam pages. The study
investigates the prevalence of malicious push advertising and the limitations of
traditional ad-blocking and URL-filtering techniques.

The authors develop PushAdMiner, which collects and analyzes web push
notifications. An instrumented Chromium browser inspects service-worker behavior,
while a crawler visits websites, registers for notifications, and collects
messages. The system groups related notifications into advertising campaigns and
examines their destinations. It collected 21,541 web push notifications and
identified 5,143 advertisements across 572 campaigns; 51\% of the collected
advertisements were reported as malicious. Traditional ad blockers and URL
filters were often unable to block them.

BrowseBlock adopts examination of notification content, notification destinations,
and suspicious notification-related indicators. These findings inform the browser
notification-scam module, which can warn users and present supporting information
through the extension. PushAdMiner was primarily a measurement and collection
framework rather than a complete real-time blocking classifier. BrowseBlock adapts
its detection insights without claiming to implement the entire framework.

\clearpage
\section{CPDetector: Automated Detection of Excessive Permissions in Chrome Extensions, Preprint}
The paper by Xiaoyu Cheng and co-authors investigates security risks associated
with extensions requesting more permissions than their functionality appears to
require. Extensions may access tabs, cookies, browsing history, and webpage
content. Excessive permissions expose sensitive information and increase the
potential for unauthorized access. The paper evaluates the relationship between
declared permissions and APIs used by an extension.

CPDetector extracts permission information from \texttt{manifest.json} and
analyzes JavaScript code to identify API calls corresponding to browser
permissions. It compares declared permissions with permissions inferred from the
implementation. PermissionGuard, a browser extension, alerts users about excessive
permissions. The framework analyzed 199,058 Chrome Web Store extensions and
identified 98,072, or 49.3\%, as requesting excessive permissions according to
its methodology. The uploaded version is a preprint and has not been peer
reviewed.

BrowseBlock adopts the principle of least privilege and permission-related risk
assessment for its browser-extension-abuse module. Accessible extension metadata,
declared permissions, and known suspicious indicators can identify potentially
risky configurations. However, a standard extension does not automatically have
unrestricted access to the internal code or runtime behavior of every other
installed extension. BrowseBlock must limit its claims to information available
through supported browser mechanisms.

\clearpage
\section{A Login Page Transparency and Visual Similarity-Based Zero-Day Phishing Defense Protocol, 2025}
The paper by Gaurav Varshney, Akanksha Raj, Divya Sangwan, Sharif Abuadbba,
Rina Mishra, and Yansong Gao proposes a phishing-detection and prevention protocol
for fraudulent login pages. It focuses on phishing websites that imitate
legitimate authentication pages and attempt to collect usernames and passwords.
The authors identify limitations of URL blacklists and heuristic detection,
which may miss newly created phishing pages. The proposed Page Transparency
concept verifies login-page legitimacy before allowing a page to be displayed.

The approach uses a Public Log Server to maintain records of legitimate login
pages and cryptographic proofs to verify registration. It also uses visual
page-matching with a Siamese Network Model to identify deceptive pages that
closely resemble legitimate ones. A custom HTTP Page Transparency header supports
client-side verification. The authors developed the AntiPhish Chrome extension
and evaluated credential-requiring pages associated with 91 brands, including
2,325 phishing pages, 115 legitimate pages, and 50 previously unseen legitimate
pages.

BrowseBlock adopts login-page identification, deceptive-authentication-interface
detection, and browser-level warnings for its credential-harvesting module. The
module can examine accessible form fields, submission destinations, URL
characteristics, and inconsistencies between displayed identity and actual domain.
Visual similarity could provide an additional signal if implemented and evaluated.
BrowseBlock does not require the complete Page Transparency infrastructure,
including its public log server and custom HTTP header. Findings can be combined
with URL analysis to warn users before they disclose sensitive credentials.
